\documentclass[11pt]{article}

\usepackage[letterpaper,margin=1in]{geometry}
\usepackage{amsmath,amssymb}
\usepackage{booktabs}
\usepackage{longtable}
\usepackage{array}
\usepackage{xcolor}
\usepackage{listings}
\usepackage{url}
\usepackage{hyperref}
\usepackage{titlesec}
\usepackage{parskip}
\usepackage{fancyhdr}
\usepackage{enumitem}

\hypersetup{
  colorlinks=true,
  linkcolor=black,
  urlcolor=blue!50!black,
  citecolor=black,
  pdftitle={Token Importance Unit --- Software Overview},
  pdfauthor={LonghornSilicon},
  pdfsubject={Compiler co-design entry point}
}

\titleformat{\section}{\normalfont\Large\bfseries}{\thesection.}{0.6em}{}
\titleformat{\subsection}{\normalfont\large\bfseries}{\thesubsection}{0.6em}{}

\pagestyle{fancy}
\fancyhf{}
\lhead{\small TIU SW Overview}
\rhead{\small Compiler Co-Design Entry Point}
\cfoot{\thepage}
\renewcommand{\headrulewidth}{0.4pt}

\definecolor{codebg}{gray}{0.97}
\definecolor{codekey}{rgb}{0.13,0.29,0.53}
\definecolor{codecom}{rgb}{0.25,0.5,0.25}
\lstset{
  basicstyle=\ttfamily\small,
  backgroundcolor=\color{codebg},
  frame=single,
  framerule=0.4pt,
  rulecolor=\color{black!30},
  xleftmargin=0.5em,
  xrightmargin=0.5em,
  showstringspaces=false,
  breaklines=true,
  keywordstyle=\color{codekey}\bfseries,
  commentstyle=\color{codecom}\itshape,
  numbers=none
}

\setlist{itemsep=2pt,topsep=3pt}

% --- table breathing room ---------------------------------------------------
% Roomier rows, a touch more column padding, and clear vertical separation
% between floats and the surrounding body text (important with parskip, where
% floats otherwise butt against paragraphs).
\renewcommand{\arraystretch}{1.25}
\setlength{\tabcolsep}{7pt}
\setlength{\textfloatsep}{18pt plus 2pt minus 2pt}
\setlength{\intextsep}{16pt plus 3pt minus 2pt}
\setlength{\floatsep}{16pt plus 2pt minus 2pt}

\begin{document}

%==============================================================================
\thispagestyle{empty}
\begin{center}
\vspace*{1.2in}
{\Huge\bfseries Token Importance Unit}\\[0.4em]
{\LARGE\bfseries Software \& Reference-Model Overview}\\[1.5em]

\rule{0.7\textwidth}{0.4pt}\\[1em]

\begin{tabular}{rl}
\textbf{Date}:       & 2026-07-18 \\
\textbf{Status}:     & Block 3 of 4 --- RTL signed off, reference model at parity \\
\textbf{Audience}:   & Compiler engineers integrating against the chip \\
\textbf{Repository}: & \texttt{LonghornSilicon/token-importance-unit} \\
\end{tabular}\\[1em]
\rule{0.7\textwidth}{0.4pt}

\vfill
\begin{minipage}{0.85\textwidth}
\small
\textbf{Purpose.} This document is the orientation map for a compiler team
integrating against the Token Importance Unit. It points to the bit-accurate
reference model, the ISA spec, the RTL and its test vectors, and the compiler
entry point. Read this first; then jump to the specific references it links.
\end{minipage}

\vspace{0.3in}
\end{center}

\clearpage
\tableofcontents
\clearpage

%==============================================================================
\section{What the TIU Is}
\label{sec:what}

The Token Importance Unit (TIU) is block~3 of the four-block LonghornSilicon LLM
inference accelerator targeting TSMC 16nm FinFET (N16FFC). It is the H2O
heavy-hitter eviction core for the KV cache: it accumulates each cached token's
received attention mass, evicts the least-attended token when the cache is full,
and emits a per-token keep/demote tier signal to the KV Cache Engine (block~2).

The silicon core is deliberately small --- \emph{accumulate + argmin + tier}. The
H2O recent-window protection and KV-budget bookkeeping are a control-layer wrapper
the compiler emits around the core, exactly as block~1 shipped just the ratio gate.

%==============================================================================
\section{The Software Stack}
\label{sec:stack}

Four artifacts make up the compiler-facing software for this block:

\begin{table}[ht]
\centering
\small
\caption{Software artifacts and their status.}
\label{tab:coverage}
\begin{tabular}{@{}p{5.4cm} p{4.0cm} p{4.4cm}@{}}
\toprule
Artifact & Role & Status \\
\midrule
\path{rtl/token_importance_unit.sv}          & The DUT (parameterized SystemVerilog) & Sky130 signed off, 0 violations \\
\path{sw/reference_model/tiu_ref.py}         & Bit-accurate Python model             & 40/40 evictions vs RTL \\
\path{rtl/tb/testvectors/tiu_trace.hex}      & Golden replay trace (real Qwen2)      & 48 tokens, 40 evictions \\
\path{sw/reference_model/example_compiler_use.py} & Compiler entry point (H2O wrapper) & Runnable \\
\bottomrule
\end{tabular}
\end{table}

\noindent The reference model is the contract: it exposes the same functional
behavior the hardware implements (same saturating accumulators, same
serialized-argmin eviction, same tier comparator), so a compiler backend can
target the TIU before silicon or an FPGA prototype exists.

%==============================================================================
\section{Where Things Live}
\label{sec:layout}

\begin{lstlisting}[basicstyle=\ttfamily\small]
rtl/
+-- token_importance_unit.sv          (the DUT)
+-- tb/
    +-- tb_token_importance_unit.sv   (directed, 30/30)
    +-- tb_realdata.sv                (real-Qwen2 replay, 40/40)
    +-- testvectors/tiu_trace.hex     (committed golden trace)

sw/reference_model/
+-- README.md                         (API reference)
+-- tiu_ref.py                        (bit-accurate model + self-test)
+-- test_tiu_ref.py                   (Python<->RTL parity on the trace)
+-- example_compiler_use.py           (budget sizing -> schedule -> tiers)

analysis/
+-- gen_tiu_testvectors.py            (regenerates the golden trace)
+-- h2o_*.py, full_stack_integration.py  (algorithm studies)

docs/
+-- isa/token_importance_unit_isa.pdf (ISA / interface spec)
+-- tier_handshake.md                 (TIU<->KVCE tier protocol)
+-- sky130_signoff.md                 (0-violation physical sign-off)
+-- sw_overview.pdf                   (this document)
\end{lstlisting}

%==============================================================================
\section{Build and Test (One Line Each)}
\label{sec:build}

\begin{lstlisting}[basicstyle=\ttfamily\small]
# RTL simulation (icarus verilog)
cd rtl && make sim            # directed TB:  30/30 pass
cd rtl && make sim_realdata   # replay TB:    40/40 evictions bit-exact

# Reference model (Python 3.10+, no dependencies)
python sw/reference_model/tiu_ref.py             # self-test
python sw/reference_model/test_tiu_ref.py        # Python<->RTL parity: 40/40
python sw/reference_model/example_compiler_use.py# compiler walkthrough
\end{lstlisting}

\noindent The parity test replays the committed golden trace through the Python
model and checks that all 40 eviction victims match the trace's expected column
(the same victims the SystemVerilog replay TB checks against the DUT). Agreement
on all 40 establishes Python$\leftrightarrow$RTL bit-exactness.

%==============================================================================
\section{The Compiler Entry Point}
\label{sec:entry}

\path{example_compiler_use.py} is the shape a compiler backend's codegen should
produce. It drives the reference model at three levels:

\begin{enumerate}
\item \textbf{Budget sizing.} Turn a context length into a cache size $C$ and
      recent window $L$ from the validated gold config (25\% KV budget, 50/50
      recency/heavy-hitter split), clamped to the physical \texttt{N\_SLOTS}.
\item \textbf{Eviction schedule.} Stream a token's attention over the cache, drive
      \texttt{LOAD}/\texttt{ACC}/\texttt{EVICT}, and collect the eviction schedule
      the backend emits to the KV cache engine. Heavy hitters are protected;
      background tokens are evicted.
\item \textbf{Keep/demote tiers.} Calibrate \texttt{tier\_threshold} from
      accumulated mass and read per-token keep (value @ CQ-8) / demote (value @
      CQ-4) tiers off the core.
\end{enumerate}

%==============================================================================
\section{Recommended Reading Order}
\label{sec:reading}

\begin{enumerate}
\item \textbf{This document} (top to bottom; you are almost done).
\item \path{docs/isa/token_importance_unit_isa.pdf}, especially \S4 (worked
      lowering example) and \S5 (tier handshake).
\item \path{docs/tier_handshake.md} for the TIU$\leftrightarrow$KVCE protocol
      (keep/demote/evict semantics), verified with the ACU in the loop.
\item \path{sw/reference_model/example_compiler_use.py} --- run it, then read it
      side by side. The shape of that script is the shape your codegen should
      produce.
\item \path{paper/token_importance_unit.pdf} for the algorithm validation, RTL
      verification, and Sky130 sign-off in full.
\end{enumerate}

%==============================================================================
\section{Integration Phases}
\label{sec:phases}

The interface (ISA $+$ reference model) is stable across all phases; only the
runtime implementation of the operations changes: Python function call
$\rightarrow$ AXI register write $\rightarrow$ PCIe MMIO.

\begin{table}[ht]
\centering
\small
\renewcommand{\arraystretch}{1.25}
\caption{Four-phase compiler integration plan.}
\label{tab:phases}
\begin{tabular}{@{}c p{3.0cm} p{4.0cm} p{4.8cm}@{}}
\toprule
Phase & Timeline & Compiler targets & We provide \\
\midrule
0 & now                     & Python reference model                    & Model, ISA, parity test \\
1 & when ZCU102/104 arrives & AXI-Lite + AXI-Stream on Zynq UltraScale+ & Vivado bitstream + PYNQ driver \\
2 & all four blocks land    & Same AXI; tier handshake to KVCE on HW    & Multi-block FPGA project \\
3 & post-tape-out (2027+)   & PCIe-attached custom chip                 & TSMC 16FFC silicon + Linux driver \\
\bottomrule
\end{tabular}
\end{table}

%==============================================================================
\section{Verification Summary}
\label{sec:verif}

\begin{table}[ht]
\centering
\small
\caption{Current verification status.}
\label{tab:verif}
\begin{tabular}{@{}l l c@{}}
\toprule
Surface & Test & Status \\
\midrule
RTL directed        & \path{tb_token_importance_unit.sv}         & 30/30 pass \\
RTL real-data replay & \path{tb_realdata.sv} (Qwen2 trace)       & 40/40 evictions \\
Python reference     & self-test (directed cases)                & pass \\
Python$\leftrightarrow$RTL parity & golden-trace replay          & 40/40 evictions \\
Physical             & Sky130 LibreLane/OpenROAD sign-off        & 0 violations \\
Three-block stack    & TIU $+$ KVCE $+$ APA on Qwen2             & $\Delta$ $-0.031$ vs FP16 \\
\bottomrule
\end{tabular}
\end{table}

\vspace{0.4in}
\noindent\rule{\textwidth}{0.4pt}\\
\noindent\small\textit{This document is the canonical entry point for a compiler
engineer working with the LonghornSilicon Token Importance Unit. It is versioned
alongside the reference model and updated as the block evolves.}

\end{document}
